% Sounio: Epistemic Types for Scientific Computing
% Target: IEEE Computing in Science & Engineering (CiSE)
%
% Keywords: uncertainty quantification, type systems, scientific computing, GUM, pharmacometrics

\documentclass[journal]{IEEEtran}

% Packages
\usepackage[utf8]{inputenc}
\usepackage[T1]{fontenc}
\usepackage{amsmath,amssymb,amsthm}
\usepackage{booktabs}
\usepackage{graphicx}
\usepackage{hyperref}
\usepackage{listings}
\usepackage{xcolor}
\usepackage{cite}
\usepackage{url}

% Code listing style
\definecolor{codebg}{RGB}{248,248,248}
\definecolor{keyword}{RGB}{0,0,180}
\definecolor{comment}{RGB}{0,128,0}
\definecolor{string}{RGB}{163,21,21}

\lstdefinelanguage{sounio}{
  keywords={fn, let, var, if, else, while, for, return, struct, enum, match, with, kernel, linear, use},
  keywordstyle=\color{keyword}\bfseries,
  ndkeywords={Knowledge, f64, i32, bool, string, IO, Alloc, GPU},
  ndkeywordstyle=\color{blue},
  sensitive=true,
  comment=[l]{//},
  commentstyle=\color{comment}\itshape,
  stringstyle=\color{string},
  morestring=[b]",
}

\lstset{
  language=sounio,
  basicstyle=\ttfamily\footnotesize,
  backgroundcolor=\color{codebg},
  frame=single,
  numbers=left,
  numberstyle=\tiny\color{gray},
  breaklines=true,
  captionpos=b,
}

% Theorems
\newtheorem{theorem}{Theorem}
\newtheorem{lemma}[theorem]{Lemma}
\newtheorem{definition}{Definition}
\newtheorem{proposition}{Proposition}

% Title
\title{Sounio: Language-Level Uncertainty Propagation for Scientific Computing}

\author{Demetrios~Chiuratto~Agourakis%
  \thanks{D.~Chiuratto~Agourakis is with the Faculdade de Ci\^{e}ncias
    M\'{e}dicas e da Sa\'{u}de, Pontif\'{i}cia Universidade Cat\'{o}lica
    de S\~{a}o Paulo (PUC-SP), Brazil, and with the Sounio Project.
    ORCID: 0009-0001-8671-8878.
    E-mail: demetrios@agourakis.med.br}%
  \and Marli~Gerenutti%
  \thanks{M.~Gerenutti is with the Faculdade de Ci\^{e}ncias
    M\'{e}dicas e da Sa\'{u}de, PUC-SP, Brazil.
    ORCID: 0000-0001-7165-646X.}%
}

\markboth{IEEE Computing in Science \& Engineering, 2026}%
{Chiuratto~Agourakis and Gerenutti: Sounio}

\begin{document}

\maketitle

\begin{abstract}
Systems programming languages lack native support for propagating
measurement uncertainty in compliance with the Guide to the Expression
of Uncertainty in Measurement (GUM). We present \textbf{Sounio}, a
systems language whose \texttt{Knowledge<T>} type carries a value, GUM
variance, confidence, and provenance, with an effect system that emits
diagnostics when uncertainty may be discarded. We also introduce
\textbf{K-AXI}, an AXI4-Stream sideband extension that streams epistemic
metadata through FPGA and GPU fabrics. Cross-validation against the
analytical Bateman equation confirms first-order convergence under
step-size refinement. Results are for first-order GUM propagation with
uncorrelated inputs. Sounio is open source (Apache~2.0).
\end{abstract}

\begin{IEEEkeywords}
Uncertainty quantification, type systems, scientific computing, GUM,
pharmacometrics, epistemic computing.
\end{IEEEkeywords}

%%% ================================================================
\section{Introduction}
\label{sec:intro}

\IEEEPARstart{S}{cientific} computing frequently involves uncertain
quantities: measured values with instrument precision, model parameters
with confidence intervals, and predictions with error bounds. Yet
mainstream programming languages treat numerical values as exact, forcing
developers to either: (1)~ignore uncertainty, compromising scientific
validity; (2)~track uncertainty manually, which is error-prone; or
(3)~use library wrappers with performance overhead and API friction.

The Guide to the Expression of Uncertainty in Measurement
(GUM)~\cite{jcgm2008gum} provides a rigorous framework for uncertainty
propagation, but implementing GUM in general-purpose languages requires
discipline that compilers cannot enforce.

We introduce \textbf{Sounio}, a systems programming language with language
features for epistemic computation, where uncertainty, confidence, and
provenance are represented as type-level metadata rather than external
library conventions. The \texttt{Knowledge<T>} type
carries not just a value but also:
\begin{itemize}
  \item \textbf{Variance}: Uncertainty in the GUM sense ($u^2$)
  \item \textbf{Confidence}: A Beta-distributed credibility channel
  \item \textbf{Provenance}: Audit trail for traceability
\end{itemize}

Arithmetic on \texttt{Knowledge<T>} automatically propagates uncertainty
using first-order Taylor series (GUM linear method). Confidence is tracked
as a separate epistemic channel and is not substituted for metrological
variance. Full Monte Carlo propagation~\cite{jcgm2008supplement1} is an
active extension and is not a claimed contribution of this manuscript.

Our contributions are:
\begin{itemize}
  \item An epistemic type system with GUM-compliant propagation and
        compile-time effect tracking (Section~\ref{sec:type-system})
  \item K-AXI, a hardware bus extension that streams epistemic metadata
        through accelerator fabrics (Section~\ref{sec:hardware})
  \item Cross-validation against analytical pharmacokinetic solutions
        with a validated dosing case study (Section~\ref{sec:evaluation})
\end{itemize}

%%% ================================================================
\section{Background}
\label{sec:background}

\subsection{The GUM Framework}

The GUM~\cite{jcgm2008gum} defines standard uncertainty $u(x)$ as the
positive square root of variance:
\[
u(x) = \sqrt{\text{Var}(x)}
\]

For a function $y = f(x_1, \ldots, x_N)$ of uncorrelated inputs, the
\emph{combined standard uncertainty} is:
\begin{equation}
  u_c^2(y) = \sum_{i=1}^{N} \left(\frac{\partial f}{\partial x_i}\right)^2 u^2(x_i)
  \label{eq:gum-propagation}
\end{equation}

This first-order Taylor expansion (the ``law of propagation of
uncertainty'') is exact for linear functions and approximate otherwise.

\subsection{Existing Approaches}

\paragraph{Python: uncertainties}
The \texttt{uncertainties} package~\cite{lebigot2024uncertainties}
provides \texttt{ufloat} objects with automatic differentiation for
(\ref{eq:gum-propagation}). It is widely used and mature, but static
guarantees about uncertainty preservation are outside Python's type system.

\paragraph{Julia: Measurements.jl}
Julia's \texttt{Measurements.jl}~\cite{giordano2016measurements} uses
dual numbers for uncertainty tracking. Related Julia tooling includes
\texttt{Uncertain.jl}~\cite{uncertainjl},
\texttt{UncertainData.jl}~\cite{uncertaindatajl}, and
\texttt{IntervalArithmetic.jl}~\cite{intervalarithmeticjl}, which offer
strong uncertainty workflows at the library level.

\paragraph{C++ Toolkits}
The Uncertainty Quantification Toolkit (UQTK)~\cite{uqtk} provides
polynomial chaos and stochastic collocation methods in C++, but does not
enforce uncertainty handling through a host-language type system.

\paragraph{Domain-Specific Languages}
Languages like Stan~\cite{carpenter2017stan} and BUGS model uncertainty
probabilistically but are not general-purpose systems languages.

\begin{table}[!t]
\caption{Representative uncertainty systems compared on key dimensions.}
\label{tab:approaches}
\centering
\footnotesize
\begin{tabular}{llllll}
\toprule
Tool & Static & GUM & Corr. & Prov. & Ecosystem \\
\midrule
uncertainties & No & Yes & Yes & No & Mature \\
Measurements.jl & No & Yes & Yes & No & Mature \\
IntervalArithmetic.jl & No & Interval & N/A & No & Mature \\
UQTK & No & PCE/MC & Yes & No & Moderate \\
\textbf{Sounio} & \textbf{Yes} & \textbf{Yes} & \textbf{No}$^{*}$ & \textbf{Yes} & \textbf{Early} \\
\bottomrule
\end{tabular}

\smallskip
\raggedright
\footnotesize
$^{*}$Correlated-input covariance propagation is planned but not yet
implemented. ``Static'' = compile-time guarantees on uncertainty preservation.
``Corr.'' = support for correlated inputs. ``Ecosystem'' maturity reflects
package availability, documentation breadth, and community adoption.
\end{table}

\subsection{Limitations of Library Approaches}

Library-based uncertainty tracking has three fundamental limitations:
\begin{enumerate}
  \item \textbf{No static guarantees}: A function accepting
        \texttt{float64} can silently drop uncertainty.
  \item \textbf{Runtime overhead}: Wrapper types require heap allocation,
        indirection, and runtime dispatch.
  \item \textbf{No provenance}: Audit trails for regulatory compliance
        must be implemented separately.
\end{enumerate}

Sounio addresses all three by making uncertainty a compile-time type
property.

%%% ================================================================
\section{Epistemic Type System and Propagation}
\label{sec:type-system}

\subsection{The Knowledge Type}

The core epistemic type is:

\begin{lstlisting}[caption={Knowledge type definition}]
struct Knowledge<T> {
    value: T,                    // Point estimate
    variance: f64,               // Uncertainty (sigma^2)
    confidence: BetaConfidence,  // Credibility
    provenance: Provenance,      // Audit trail
}
\end{lstlisting}

\texttt{BetaConfidence} represents epistemic confidence as a Beta
distribution, supporting Bayesian updating:

\begin{lstlisting}[caption={Confidence representation}]
struct BetaConfidence {
    alpha: f64,  // Successes + prior
    beta: f64,   // Failures + prior
}

fn beta_mean(c: &BetaConfidence) -> f64 {
    c.alpha / (c.alpha + c.beta)
}
\end{lstlisting}

\subsection{Compiler Warnings for Uncertainty Loss}

Sounio's type system prevents silent uncertainty loss through
compile-time warnings. When a \texttt{Knowledge<T>} value is extracted to
\texttt{T}, the compiler issues a diagnostic:

\begin{lstlisting}[caption={Illustrative extraction diagnostic (pseudocode)}]
fn main() with IO {
    let k = Knowledge::measured(5.0, 0.01, "scale")
    let bad = compute_exact(extract_with_warning(k))
    // warning: uncertainty discarded by extract_with_warning
    let good = k.scale(2.0)
    // OK: uncertainty preserved
}
\end{lstlisting}

\begin{definition}[Epistemic Extraction]
The operation $\text{extract}: \texttt{Knowledge<T>} \to \texttt{T}$ is
explicit and emits a compiler diagnostic at the call site to signal that
uncertainty information is being discarded. Implicit coercion from
\texttt{Knowledge<T>} to \texttt{T} is disallowed.
\end{definition}

\subsection{Provenance Tracking}

Every \texttt{Knowledge} value carries a \texttt{Provenance} chain:

\begin{lstlisting}[caption={Provenance structure}]
enum Source {
    Measurement { instrument: String, timestamp: i64 },
    Computed { operation: String },
    Assertion { author: String },
    External { url: String, doi: Option<String> },
}

struct Provenance {
    sources: Vec<Source>,
    merkle_hash: [u8; 32],
}
\end{lstlisting}

This provides structural support for audit trails relevant to regulated
computational systems (e.g., FDA 21 CFR Part~11~\cite{fda21cfr11},
ISO~17025~\cite{iso17025}).

\subsection{Formal Properties}

The type system enforces two key invariants: (1)~\texttt{Knowledge<T>}
cannot implicitly coerce to~\texttt{T}, and (2)~arithmetic on epistemic
values produces epistemic results with variance computed by the GUM law
of propagation (\ref{eq:gum-propagation}). Progress, preservation, and
first-order GUM soundness for the core epistemic calculus are established
in a mechanized Lean~4 artifact (\texttt{formal/lean4/}, 728
definitions/theorems, zero \texttt{sorry}).

\subsection{First-Order (Linear) Propagation}
\label{sec:propagation}

For arithmetic operations, Sounio implements (\ref{eq:gum-propagation})
at the IR level. The compiler transforms:
\begin{lstlisting}
let z = x * y  // x, y: Knowledge<f64>
\end{lstlisting}
into:
\begin{lstlisting}
let z_val = x.value * y.value
let z_var = y.value^2 * x.variance
          + x.value^2 * y.variance
let z = Knowledge { value: z_val,
                    variance: z_var, ... }
\end{lstlisting}

\subsection{Propagation Correctness}

\begin{theorem}[First-Order GUM Conformance]
\label{thm:gum-compliance}
For any expression $e$ composed of arithmetic operations on
\texttt{Knowledge} values, let
$\text{eval}(e) = (v, \sigma^2, \beta, \pi)$ under Sounio's operational
semantics. Then $\sigma^2 = u_c^2(e)$ as defined by
(\ref{eq:gum-propagation}), under the assumptions of first-order
linearization and uncorrelated inputs.
\end{theorem}

\begin{proof}
By structural induction on $e$.

\textbf{Base case:}
$e = \texttt{Knowledge::new}(v, \sigma^2, \beta, \pi)$.
By definition, $\text{eval}(e) = (v, \sigma^2, \beta, \pi)$ and
$u_c^2(e) = \sigma^2$. Thus $\sigma^2 = u_c^2(e)$.

\textbf{Inductive case (addition):} Let $e = e_1 + e_2$ where by IH,
$\text{eval}(e_i) = (v_i, \sigma_i^2, \ldots)$ with
$\sigma_i^2 = u_c^2(e_i)$. By E-ADD,
$\text{eval}(e) = (v_1{+}v_2, \sigma_1^2{+}\sigma_2^2, \ldots)$.
For $f(x,y) = x+y$, GUM gives
$u_c^2 = (\partial f/\partial x)^2 \sigma_1^2 +
(\partial f/\partial y)^2 \sigma_2^2 = \sigma_1^2 + \sigma_2^2$.
\checkmark

\textbf{Inductive case (multiplication):} Let $e = e_1 \times e_2$.
By E-MUL,
$\text{eval}(e) = (v_1 v_2, v_2^2\sigma_1^2 + v_1^2\sigma_2^2, \ldots)$.
For $f(x,y) = xy$, $\partial f/\partial x = y$ and
$\partial f/\partial y = x$, so
$u_c^2 = v_2^2\sigma_1^2 + v_1^2\sigma_2^2$. \checkmark

Division and transcendentals follow similarly via the delta method.
\end{proof}

\subsection{Transcendental Functions}

For $y = f(x)$, the delta method gives:
\[
  u^2(y) \approx \left(\frac{df}{dx}\right)^2 u^2(x)
\]

Sounio provides built-in derivatives for standard functions
(Table~\ref{tab:transcendental}).

\begin{table}[!t]
\caption{Representative built-in uncertainty propagation rules (delta method).}
\label{tab:transcendental}
\centering
\begin{tabular}{lll}
\toprule
Function & $f(x)$ & $u^2(f(x))$ \\
\midrule
\texttt{exp} & $e^x$ & $e^{2x} u^2(x)$ \\
\texttt{ln} & $\ln x$ & $u^2(x)/x^2$ \\
\texttt{sin} & $\sin x$ & $\cos^2(x)\, u^2(x)$ \\
\texttt{inv} & $1/x$ & $u^2(x)/x^4$ \\
\texttt{pow(a)} & $x^a$ & $a^2x^{2a-2}u^2(x)$ \\
\bottomrule
\end{tabular}
\end{table}

The compiler provides analogous rules for all standard transcendentals
(sqrt, cos, tan, asin, acos, atan, log10). Singularities and branch cuts
are handled with runtime domain checks.

\subsection{Scientific IR (SIR) Design}

Sounio's compiler includes a Scientific IR that represents epistemic
operations as first-class nodes:

\begin{lstlisting}[caption={SIR epistemic instructions}]
enum SirInst {
    PropagateAdd { lhs: NodeId, rhs: NodeId,
                   result: NodeId },
    PropagateMul { lhs: NodeId, rhs: NodeId,
                   result: NodeId },
    PropagateDiv { lhs: NodeId, rhs: NodeId,
                   result: NodeId },
    PropagatePow { base: NodeId, exponent: f64,
                   result: NodeId },
    PropagateTranscendental { arg: NodeId, fn_id: FnId,
                              result: NodeId },
}
\end{lstlisting}

The SIR enables domain-specific optimizations: \emph{variance fusion}
(fusing $u^2(x + y + z)$ into a single pass), \emph{exact elision}
(skipping variance computation for zero-variance operands), and
\emph{provenance elision} (skipping Merkle hash computation when
provenance is not required).

%%% ================================================================
\section{Implementation}
\label{sec:implementation}

\subsection{Compiler Architecture}

Sounio's compilation pipeline is organized as:
\[
\text{Source} \!\to\! \text{Lexer} \!\to\! \text{Parser} \!\to\! \text{AST} \!\to\!
\text{Check} \!\to\! \text{HIR} \!\to\! \text{SIR} \!\to\! \text{HLIR} \!\to\! \text{Backend}
\]

The pipeline includes four code generation backends (native ELF/Mach-O,
Cranelift JIT, partial LLVM, and GPU via PTX/SPIR-V). Developer tooling
(LSP server, formatter, linter, REPL) is implemented in the same
repository.

The compiler is fully self-hosted: a 138{,}000-line compiler and
272{,}000-line scientific standard library are written entirely in Sounio
itself. A three-stage verified bootstrap ensures correctness: the Rust-hosted
Stage~1 compiler produces a Sounio binary (Stage~2), which then compiles
itself (Stage~3); SHA-256 parity between Stage~2 and Stage~3 outputs
verifies deterministic compilation. The full repository totals approximately
690{,}000 lines of Sounio code including tests and developer tooling.

\subsection{Memory Layout}

\texttt{Knowledge<f64>} occupies 40 bytes, stack-allocated
(Table~\ref{tab:layout}). The layout is designed for SIMD-friendly
access; the native backend uses SIMD (AVX2, AVX-512, NEON) for
quaternion and octonion algebra.

\begin{table}[!t]
\caption{Memory layout of \texttt{Knowledge<f64>}.}
\label{tab:layout}
\centering
\begin{tabular}{lr}
\toprule
Field & Size \\
\midrule
\texttt{value} & 8 bytes (f64) \\
\texttt{variance} & 8 bytes (f64) \\
\texttt{confidence} & 16 bytes (BetaConf) \\
\texttt{provenance} & 8 bytes (pointer) \\
\midrule
\textbf{Total} & \textbf{40 bytes} \\
\bottomrule
\end{tabular}
\end{table}

%%% ================================================================
\section{Hardware-Aware Epistemic Propagation}
\label{sec:hardware}

Software-only uncertainty propagation encounters a fundamental limitation
when computations move to accelerator fabrics: standard bus protocols
(AXI4, PCIe, NVLink) carry data payloads but discard epistemic metadata.
A GPU kernel that receives a \texttt{Knowledge<f64>} value through a
conventional DMA transfer must reconstruct variance, confidence, and
provenance from separate memory transactions, breaking the atomicity
guarantee that the type system provides on the CPU side.

Sounio addresses this with \textbf{K-AXI} (Knowledge AXI), a bus protocol
extension that carries epistemic sideband information alongside data in
every transaction. K-AXI preserves the invariant that no epistemic metadata
is silently dropped when crossing a hardware boundary.

\subsection{K-AXI Bus Protocol}

K-AXI extends the AXI4-Stream \texttt{TUSER} sideband to carry a 480-bit
epistemic payload alongside each 64-bit data word. The sideband is
structured as:

\begin{table}[!t]
\caption{K-AXI epistemic sideband layout (480 bits in \texttt{TUSER}).}
\label{tab:kaxi-sideband}
\centering
\footnotesize
\begin{tabular}{lrl}
\toprule
Field & Width & Encoding \\
\midrule
\texttt{variance} & 64 bits & Q0.64 fixed-point \\
\texttt{provenance} & 256 bits & Merkle256 digest \\
\texttt{confidence} & 128 bits & Beta$(\alpha, \beta)$, 2$\times$64-bit \\
\texttt{flags} & 32 bits & \{$\alpha_{q8}$, succ, fail, op\} \\
\midrule
\textbf{Total sideband} & \textbf{480 bits} & \\
\bottomrule
\end{tabular}
\end{table}

The \texttt{flags} field encodes the operation kind (ADD, MUL, DIV, FMA),
8-bit Bayesian update deltas (\texttt{success\_inc}, \texttt{failure\_inc}),
and an epistemic degradation coefficient $\alpha_{q8}$ in Q0.8 format.
A master module packs \texttt{Knowledge<f64>} values into this format; a
slave module unpacks, propagates variance, updates confidence, and folds
provenance in a single clock cycle.

\subsection{Variance Propagation in Hardware}

The K-AXI slave implements first-order variance growth using bit-shift
approximations of the delta method. For each operation kind, variance is
updated as:

\begin{table}[!t]
\caption{Hardware variance growth factors by operation kind.}
\label{tab:hw-variance}
\centering
\begin{tabular}{llr}
\toprule
Operation & Shift expression & Growth \\
\midrule
ADD & $\sigma^2 + (\sigma^2 \gg 3)$ & $+12.5\%$ \\
MUL & $\sigma^2 + (\sigma^2 \gg 1) + (\sigma^2 \gg 3)$ & $+62.5\%$ \\
DIV & $\sigma^2 + (\sigma^2 \gg 1) + (\sigma^2 \gg 2)$ & $+75.0\%$ \\
FMA & $\sigma^2 + (\sigma^2 \gg 1) + (\sigma^2 \gg 2) + (\sigma^2 \gg 4)$ & $+81.25\%$ \\
\bottomrule
\end{tabular}
\end{table}

These are conservative upper bounds; the shift-based implementation avoids
floating-point multipliers entirely, using only 64-bit adders and barrel
shifters. An additional epistemic degradation term
$\Delta\sigma^2 = \sigma^2_{\text{op}} \cdot \alpha_{q8} / 256$ models
information loss from quantization or approximation.

\subsection{Hardware Provenance and Confidence}

Provenance is maintained as a 256-bit Merkle digest, folded
deterministically via XOR rotation with a compile-time salt:
\[
  \pi' = \text{rot}(\pi, 8) \oplus \text{pad}(v) \oplus \text{pad}(\textit{flags}) \oplus S_{\text{K-AXI}}
\]
where $S_{\text{K-AXI}}$ is a fixed 256-bit salt derived from the protocol
version identifier. An optional external Merkle lane supports
cryptographic hash cores when available; the internal XOR fold serves as a
deterministic fallback.

Confidence is encoded as a Beta$(\alpha, \beta)$ pair in 128 bits
(2$\times$64-bit integers). Each transaction atomically increments $\alpha$
and $\beta$ by the 8-bit deltas in the flags field, supporting online
Bayesian updating of measurement credibility.

\subsection{Epistemic Power Accumulator}

The \textbf{Epistemic Power Accumulator} is a hardware module that
aggregates counter increments from the epistemic fabric and produces a
live Q32.32 fixed-point score:

\begin{equation}
  \mathcal{E}_{\text{hw}} = \frac{F}{8} + \frac{P}{16}
  \label{eq:epistemic-power}
\end{equation}

\noindent where $F$ is the accumulated fidelity (GUM-compliant transaction
count) and $P$ is the provenance depth (number of operations in the
Merkle chain). The weighting reflects the relative contribution of each
epistemic channel: runtime fidelity dominates ($\times 1/8$) while
provenance provides secondary assurance ($\times 1/16$).
Division is implemented as right-shifts, avoiding floating-point logic
entirely. In practice, the host runtime reads $\mathcal{E}_{\text{hw}}$
via a memory-mapped register; a drop below a configurable threshold
signals that accumulated epistemic degradation has exceeded tolerance,
triggering recalibration or alerting the operator.

\subsection{GPU Integration}

On the GPU side, K-AXI is realized through a 64-entry atomic ring buffer
in device memory. When a Sounio \texttt{kernel} function operates on
\texttt{Knowledge<T>} values, the compiler emits calls to a CUDA runtime
intrinsic (shown in C for clarity):

\begin{lstlisting}[caption={GPU-side K-AXI publish (compiler-emitted CUDA runtime intrinsic, not user-written Sounio code)},language=C]
__device__ void __knowledge_kaxi_publish(
    double value, double variance,
    uint256_t prov, double confidence,
    uint32_t flags)
{
    unsigned idx = atomicAdd(&g_kaxi_write_idx, 1u)
                 & (KAXI_RING_SIZE - 1);
    g_kaxi_ring[idx] = {value, variance,
        prov, __double_as_longlong(confidence),
        flags, 0};
}
\end{lstlisting}

This enables epistemic metadata to flow from GPU kernels back to the
host runtime or to downstream FPGA consumers without breaking the
\texttt{Knowledge<T>} invariant across the heterogeneous boundary.

\subsection{Synthesis and Verification}

The K-AXI modules are synthesized with Yosys targeting generic gate-level
netlists. The complete epistemic slave (variance propagation, confidence
update, Merkle fold) synthesizes to approximately 50K~equivalent gates in
a single combinatorial stage. The epistemic MAC (fused multiply-accumulate
with delta-method variance) requires two 64$\times$64 signed multipliers
and synthesizes to approximately 200K~gates.

Four testbenches validate the hardware implementation:
(1)~master-slave handshake with three operation profiles (ADD, FMA, Merkle);
(2)~bidirectional FIFO overflow and drain;
(3)~epistemic power accumulation with clear-and-resume; and
(4)~delta-method MAC with signed operands.
A 544-bit return packet format carries telemetry (timestamp, digest,
counters, flags) from the fabric to the host for runtime inspection.

%%% ================================================================
\section{Evaluation}
\label{sec:evaluation}

\subsection{Performance Characteristics}

Sounio's design offers structural properties that reduce avoidable runtime
overhead in uncertainty propagation:
\begin{enumerate}
  \item \textbf{Stack allocation}: \texttt{Knowledge<f64>} is a 40-byte
        stack-allocated value.
  \item \textbf{Static dispatch}: All epistemic operations are
        monomorphized at compile time, eliminating runtime dispatch
        overhead.
  \item \textbf{Compiler-visible structure}: Because uncertainty
        propagation is part of the IR, the compiler can apply
        domain-specific optimizations.
\end{enumerate}

Cross-language benchmarks compare Python's \texttt{uncertainties},
Julia's \texttt{Measurements.jl}, and Sounio's \texttt{Knowledge<T>}
on identical workloads. All benchmarks use the same operator mix
(addition, subtraction, multiplication, division) and problem sizes.

\textbf{Important caveat:} Sounio timings are from the tree-walking
interpreter; Julia uses LLVM JIT compilation. Raw throughput comparisons
between an interpreter and a JIT compiler are not meaningful.
The relevant finding is \emph{asymptotic scaling behavior} under
chain accumulation, where algorithmic differences dominate constant
factors. Table~\ref{tab:uncertainty-benchmarks} reports median
wall-clock times over 10 runs after warmup on an Intel Xeon Gold
6148 (2.40\,GHz, 20 cores, 128\,GB RAM, Ubuntu 22.04).

\begin{table}[!t]
\caption{Uncertainty propagation benchmarks (median ms). ``DNF'' = did not finish within 600\,s.}
\label{tab:uncertainty-benchmarks}
\centering
\footnotesize
\begin{tabular}{lrrr}
\toprule
Benchmark & Python & Julia & Sounio \\
 & \texttt{uncertainties} & \texttt{Measurements.jl} & \texttt{Knowledge<T>} \\
\midrule
GUM Chain (10K ops) & 137 & DNF$^{\dagger}$ & \textbf{29} \\
GUM Chain (100K ops) & 1\,804 & DNF$^{\dagger}$ & \textbf{237} \\
Monte Carlo (10K) & 652 & \textbf{5} & 51 \\
Monte Carlo (100K) & 6\,496 & \textbf{49} & 455 \\
Cockcroft--Gault (10K) & 606 & DNF$^{\dagger}$ & \textbf{71} \\
Transcendental (10K) & 584 & \textbf{4}$^{\ddagger}$ & 53 \\
\bottomrule
\end{tabular}
\smallskip

\raggedright
\footnotesize
$^{\dagger}$Measurements.jl accumulates a derivative dictionary per unique
measurement source; long accumulating chains exhibit $O(n^2)$ runtime
(GUM Chain 1K: 151\,ms, extrapolated 10K: $>$\,60\,s).
Independent-per-iteration workloads (Monte Carlo) remain $O(n)$.

$^{\ddagger}$Extrapolated from 1K benchmark (0.37\,ms).
\end{table}

\paragraph{Scaling behavior}
The key finding is not raw speed but \emph{asymptotic scaling under
chain accumulation}. Measurements.jl's derivative-tracking approach
is~$O(1)$ per operation for short chains but degrades to $O(n)$ per
operation when a single measurement accumulates derivatives across
many steps (Cockcroft--Gault at 1K patients: 7.6\,s; extrapolated
10K: $>$\,600\,s). Python's \texttt{uncertainties} uses a similar
derivative-tracking strategy but with CPython overhead, resulting in
consistent $O(n)$ scaling. Sounio's \texttt{Knowledge<T>} propagates
variance in $O(1)$ per operation regardless of chain length,
achieving linear scaling on all workloads.

\subsection{Numerical Accuracy: ODE Solver Cross-Validation}

To validate Sounio's numerical computing infrastructure, we cross-validate
against the analytical Bateman equation for a two-compartment
pharmacokinetic model~\cite{gibaldi1982pharmacokinetics}. This tests
floating-point arithmetic accuracy and long-running loop execution with
struct mutation (up to 2,400 iterations), exercising the compiler's SSA
construction, phi-node insertion, and backend codegen.

\paragraph{Model}
The two-compartment oral absorption model uses parameters:
Dose $= 500\,\text{mg}$, absorption rate $k_a = 1.0\,\text{h}^{-1}$,
elimination rate $k_e = 0.1\,\text{h}^{-1}$.
The closed-form Bateman solution gives:
\begin{align}
  C_{\text{gut}}(t) &= D \cdot e^{-k_a t} \\
  C_{\text{central}}(t) &= \frac{D \cdot k_a}{k_a - k_e}
    \left(e^{-k_e t} - e^{-k_a t}\right)
\end{align}

\paragraph{Results}
The test suite validates Sounio's forward Euler solver at two step sizes
($\Delta t = 0.1\,\text{h}$ and $\Delta t = 0.01\,\text{h}$) against the
analytical Bateman solution. All intermediate checkpoints at
$t \in \{2, 6, 12, 24\}\,\text{h}$ fall within $5\%$ of the analytical
values. Table~\ref{tab:bateman-errors} reports central-compartment relative
errors from the reproducible reference script
\texttt{paper/scripts/bateman\_convergence.py}. Reducing $\Delta t$ from
0.1 to 0.01 reduces error by approximately one order of magnitude, as
expected for first-order global Euler error.

\begin{table}[!t]
\caption{Relative error (\%) vs analytical Bateman solution (central compartment).}
\label{tab:bateman-errors}
\centering
\begin{tabular}{lrr}
\toprule
Time point & $\Delta t=0.1$ h & $\Delta t=0.01$ h \\
\midrule
$t=2$ h & 1.893 & 0.186 \\
$t=6$ h & 0.178 & 0.017 \\
$t=12$ h & 0.601 & 0.060 \\
$t=24$ h & 1.201 & 0.120 \\
\bottomrule
\end{tabular}
\end{table}

The Euler solver ($\Delta t = 0.1\,\text{h}$, 240 steps) recovers the
expected PK profile: $C_{\max} \approx 387\,\text{mg}$ and
$\text{AUC}_{0 \to 24} \approx 5000\,\text{mg}\!\cdot\!\text{h}$
(theoretical $\text{AUC}_{0 \to \infty} = D/k_e = 5000$).

\paragraph{Higher-Order Integration}
A three-compartment model (Gut $\to$ Liver $\to$ Central) with hepatic
first-pass metabolism ($f_a = 0.8$, $Q_h = 0.5$, $\text{CL}_h = 0.3$)
was validated using the RK2 midpoint method (480 steps,
$\Delta t = 0.05\,\text{h}$). Conservation-of-mass invariants hold and
pharmacokinetic endpoints ($C_{\max}$, $T_{\max}$, AUC) fall within
physiologically plausible ranges verified by automated tests.

\paragraph{Scope and Current Limitations}
The evaluation in this paper covers first-order propagation with
uncorrelated inputs. Correlated-input covariance propagation and full GUM
Supplement~1 Monte Carlo execution inside SIR are planned extensions and
are not claimed as implemented in this manuscript.

\subsection{Reproducibility Artifacts}
The repository includes a reproducibility entrypoint:
\texttt{paper/reproduce.sh}. It builds the IEEE manuscript, executes the
cross-language baseline benchmark harness (\texttt{benchmarks/comparison}),
generates Bateman convergence CSV artifacts, and checks bibliography
references via \texttt{paper/scripts/check\_references.py}. This is
intended to keep claims in sync with repository state at submission time.

\subsection{Application: Pharmacokinetic Dosing}
\label{sec:applications}

We illustrate Sounio's language design with a validated pharmacokinetic
example. The Cockcroft--Gault equation~\cite{cockcroft1976prediction}
for creatinine clearance is widely used in clinical drug dosing:
\[
  \text{CrCl} = \frac{(140 - \text{age}) \times \text{weight} \times [0.85 \text{ if female}]}{72 \times \text{SCr}}
\]

In clinical practice, each input carries measurement uncertainty.
Sounio automatically propagates these through the formula:

\begin{lstlisting}[caption={GUM-compliant drug dosing}]
fn creatinine_clearance(
    age: Knowledge<years>,
    weight: Knowledge<kg>,
    scr: Knowledge<mg_per_dL>,
    is_female: bool
) -> Knowledge<mL_per_min> {
    let factor = if is_female { 0.85 }
                 else { 1.0 }
    let num = (Knowledge::exact_with_provenance(140.0, "Cockcroft-Gault constant") - age)
              * weight
    num.scale(factor) /
        (Knowledge::exact_with_provenance(72.0, "Cockcroft-Gault constant") * scr)
}
\end{lstlisting}

The result includes propagated uncertainty with full provenance metadata.
This example is validated by the Bateman ODE cross-validation tests in
Section~\ref{sec:evaluation}, which exercise the same uncertainty
propagation pipeline on a two-compartment model.

Sounio's standard library also includes DARWIN, a 6,000-line PBPK
modeling framework implementing 14-compartment mammalian physiology
with Rodgers--Rowland tissue partitioning~\cite{rodgers2005tissue}.

%%% ================================================================
\section{Related Work}
\label{sec:related}

\paragraph{Probabilistic Programming Languages}
Stan~\cite{carpenter2017stan} employs HMC for Bayesian inference.
Pyro~\cite{bingham2019pyro} extends deep probabilistic programming to
PyTorch, while Gen~\cite{cusumano2019gen} introduces programmable
inference. Turing.jl~\cite{ge2018turing} provides universal probabilistic
programming in Julia. These are domain-specific languages for Bayesian
inference rather than systems programming languages with native
uncertainty types. Sounio's \texttt{Knowledge<T>} tracks frequentist
(GUM) uncertainty at the type level, complementing rather than replacing
Bayesian approaches. In particular, Stan's \texttt{generated quantities}
workflow supports posterior predictive reporting, while Sounio targets
deterministic, type-level uncertainty propagation in compiled systems code.

\paragraph{Uncertainty Propagation Libraries}
Python's \texttt{uncertainties}~\cite{lebigot2024uncertainties} uses
operator overloading but incurs runtime overhead. Julia's
\texttt{Measurements.jl}~\cite{giordano2016measurements} leverages dual
numbers with JIT compilation. Both lack compile-time guarantees on
uncertainty preservation. Measure-theoretic and interval-focused Julia
ecosystems (e.g., \texttt{MeasureTheory.jl}~\cite{measuretheoryjl} and
\texttt{IntervalArithmetic.jl}~\cite{intervalarithmeticjl}) provide strong
modeling abstractions and validated bounds, but are library-level rather
than language-level typing disciplines.

\paragraph{Dependent and Refinement Types}
Idris~\cite{brady2013idris} supports full dependent types. F*~\cite{swamy2016dependent}
uses refinement types for verification. These systems model correctness
properties but not epistemic uncertainty semantically.

\paragraph{Units of Measure}
F\#'s units of measure~\cite{kennedy2009units} provide compile-time
dimensional analysis. Sounio extends this by combining units with
epistemic uncertainty.

\paragraph{Automatic Differentiation}
JAX~\cite{bradbury2018jax} and PyTorch~\cite{paszke2019pytorch} compute
gradients via AD. Sounio adapts similar techniques for uncertainty
propagation as a first-class type-level concern.

\paragraph{Linear Types and Algebraic Effects}
Wadler's linear types~\cite{wadler1990linear} and Plotkin and Pretnar's
algebraic effect handlers~\cite{plotkin2009handlers} inform Sounio's
resource management and effect system, respectively.

%%% ================================================================
\section{Conclusion}
\label{sec:conclusion}

Sounio demonstrates that uncertainty quantification can be integrated
into a systems programming language with explicit type-level metadata.
The \texttt{Knowledge<T>} type supports first-order GUM propagation for
uncorrelated inputs, while the algebraic effect system flags explicit
uncertainty erasure points. Cross-validation against analytical
pharmacokinetic solutions supports first-order convergence behavior in ODE
integration workflows.

The standard library provides extensions beyond the core type system,
including forward-mode AD with epistemic metadata, PBPK modeling tools,
and causal-identifiability checks via SMT solving.

\paragraph{Future Work}
Priority extensions are: (1)~full Monte Carlo execution in SIR;
(2)~correlated-input covariance propagation; (3)~expanded mechanized
proof coverage; and (4)~external compliance audit support for regulated
settings (e.g., FDA 21 CFR Part 11).

\paragraph{Availability}
Sounio is open source under the Apache License~2.0.
Repository: \url{https://github.com/sounio-lang/sounio}.
Documentation: \url{https://souniolang.org}.
DOI: \href{https://doi.org/10.5281/zenodo.18404188}{10.5281/zenodo.18404188}.

\section*{Acknowledgments}
The authors thank early users who reviewed language behavior,
diagnostics, and documentation during pre-release iterations.

\bibliographystyle{IEEEtran}
\bibliography{references}

\end{document}
